(1) 一维均匀流体沿 $+x$ 方向以恒速 $U\ge0$ 流动，流体相对传播速度为 $c>0$。小扰动满足 $(\partial_t+U\partial_x)^2\eta=c^2\eta_{xx}$，求平面波色散关系。
(2) 求两支波在实验室坐标系中的相速度、群速度以及流体共动系中的传播速度。
(3) 对实验室固定角频率 $\omega>0$，求两支波数，分析 $U<c$、$U=c$、$U>c$ 时能否有向上游传播的信号。
(4) 检查 $U\to0$ 的左右对称极限，并说明均匀流模型中“地平线”的确切动力学含义及其局限。
